\documentclass[12pt,a4paper]{article}

% ─────────────────────────────────────────────────────────────
% PACKAGES ET CONFIGURATION
% ─────────────────────────────────────────────────────────────
\usepackage[utf8]{inputenc}
\usepackage[T1]{fontenc}
\usepackage[french]{babel}
\usepackage{lmodern}
\usepackage{geometry}
\geometry{top=2.5cm, bottom=2.5cm, left=2.5cm, right=2.5cm}

\usepackage{graphicx}
\usepackage{xcolor}
\usepackage{listings}
\usepackage{tcolorbox}
\tcbuselibrary{skins,breakable}
\usepackage{booktabs}
\usepackage{tabularx}
\usepackage{hyperref}
\usepackage{fancyhdr}
\usepackage{titlesec}
\usepackage{float}
\usepackage{caption}
\usepackage{amsmath,amssymb}

% Couleurs personnalisées
\definecolor{primaryblue}{RGB}{24, 76, 120}
\definecolor{secondaryblue}{RGB}{41, 128, 185}
\definecolor{accentgreen}{RGB}{39, 174, 96}
\definecolor{accentred}{RGB}{192, 57, 43}
\definecolor{lightgray}{gray}{0.96}
\definecolor{codegray}{gray}{0.92}
\definecolor{darkgray}{gray}{0.3}

% Configuration des hyperliens
\hypersetup{
    colorlinks=true,
    linkcolor=primaryblue,
    citecolor=secondaryblue,
    urlcolor=secondaryblue,
    pdftitle={Rapport Shift-Right Testing - StudentManagement},
    pdfauthor={Équipe Software Testing}
}

% Configuration du style des pages
\pagestyle{fancy}
\fancyhf{}
\fancyhead[L]{\textcolor{darkgray}{\small Shift-Right Testing : Déploiement Canary \& Feature Flags}}
\fancyhead[R]{\textcolor{darkgray}{\small \thepage}}
\renewcommand{\headrulewidth}{0.5pt}
\renewcommand{\footrulewidth}{0.5pt}
\fancyfoot[C]{\textcolor{darkgray}{\footnotesize École Mohammadia d'Ingénieurs — Module Software Testing}}

% Style des titres
\titleformat{\section}
  {\color{primaryblue}\normalfont\Large\bfseries}
  {\color{primaryblue}\thesection}{1em}{}
\titleformat{\subsection}
  {\color{secondaryblue}\normalfont\large\bfseries}
  {\color{secondaryblue}\thesubsection}{1em}{}
\titleformat{\subsubsection}
  {\color{darkgray}\normalfont\normalsize\bfseries}
  {\color{darkgray}\thesubsubsection}{1em}{}

% Configuration listings pour le code
\lstset{
    backgroundcolor=\color{lightgray},
    basicstyle=\ttfamily\footnotesize,
    breaklines=true,
    captionpos=b,
    commentstyle=\color{accentgreen},
    keywordstyle=\color{primaryblue}\bfseries,
    stringstyle=\color{accentred},
    numberstyle=\tiny\color{darkgray},
    numbers=left,
    numbersep=8pt,
    showstringspaces=false,
    frame=lines,
    rulecolor=\color{secondaryblue},
    tabsize=2
}

% Commande pour insérer un placeholder de capture d'écran
\newcommand{\screenshotplaceholder}[2][7cm]{%
  \begin{figure}[H]
    \centering
    \begin{tcolorbox}[
      width=0.95\textwidth,
      height=#1,
      valign=center,
      halign=center,
      colback=lightgray,
      colframe=secondaryblue!70,
      boxrule=1.5pt,
      dashed
    ]
      \centering
      \vspace{0.2cm}
      {\color{secondaryblue}\Huge\textbf{[\,CAPTURE D'ÉCRAN\,]}}\\[0.4cm]
      {\color{darkgray}\large\textbf{#2}}\\[0.2cm]
      {\small\textit{Emplacement réservé : insérez votre capture d'écran ici}}
      \vspace{0.2cm}
    \end{tcolorbox}
    \caption{#2}
    \label{fig:#2}
  \end{figure}
}

% ─────────────────────────────────────────────────────────────
% PAGE DE GARDE
% ─────────────────────────────────────────────────────────────
\begin{document}

\begin{titlepage}
    \centering
    \vspace*{0.5cm}
    {\scshape\LARGE École Mohammadia d'Ingénieurs \par}
    \vspace{0.3cm}
    {\scshape\Large Département Génie Informatique \& Mathématiques Appliquées\par}
    \vspace{1.5cm}

    \rule{\linewidth}{0.6mm} \\[0.4cm]
    {\huge\bfseries\color{primaryblue} RAPPORT DE PROJET \par}
    \vspace{0.3cm}
    {\LARGE\bfseries Shift-Right Testing en Production : Déploiement Canary, Feature Flags \& Rollback Automatique \par}
    \rule{\linewidth}{0.6mm} \\[1.5cm]

    \textbf{\Large Projet StudentManagement}\\[0.3cm]
    {\large Spring Boot 4.1.1 — Angular 19 — Nginx — Unleash — k6 — AWS EC2}\\[2.5cm]

    \begin{minipage}{0.48\textwidth}
        \begin{flushleft} \large
            \emph{Réalisé par :}\\
            \textbf{Équipe de Test Logiciel}
        \end{flushleft}
    \end{minipage}
    ~
    \begin{minipage}{0.48\textwidth}
        \begin{flushright} \large
            \emph{Module :}\\
            \textbf{Software Testing \& Qualité}\\
            \emph{Semestre 5 (2026/2027)}
        \end{flushright}
    \end{minipage}

    \vfill
    {\large \today \par}
\end{titlepage}

\tableofcontents
\newpage

% ─────────────────────────────────────────────────────────────
% 1. INTRODUCTION & CONTEXTE
% ─────────────────────────────────────────────────────────────
\section{Introduction et Contexte du Projet}

\subsection{Problématique du Shift-Right Testing}
Traditionnellement, l'assurance qualité logicielle se concentre sur l'approche dite \textit{Shift-Left}, qui consiste à multiplier les tests au plus tôt dans le cycle de développement (tests unitaires, tests d'intégration, analyse statique). Bien qu'indispensable, cette approche ne garantit pas à elle seule l'absence de régressions ou de comportements inattendus lors de la mise en production réelle.

Le \textbf{Shift-Right Testing} complète ce paradigme en étendant les pratiques de vérification, de test de charge, d'injection de pannes et d'observabilité directement sur l'environnement de production ou pré-production, au plus près des utilisateurs réels.

\subsection{Contexte Applicatif et Contraintes Techniques}
Le projet support de cette mise en œuvre est l'application \textbf{StudentManagement}, une plateforme de gestion académique composée d'une API REST Spring Boot (Java 21) et d'un frontend en Angular 19, conteneurisés via Docker Compose.

Le cahier des charges impose la mise en œuvre d'une architecture résiliente sans surcoût d'infrastructure :
\begin{itemize}
    \item \textbf{Cible de déploiement} : Instance cloud AWS EC2 (gabarit restreint \texttt{t2.micro} / \texttt{t3.micro}, 1 vCPU, 1 Go RAM).
    \item \textbf{Contrainte d'outillage} : Déploiement sans orchestration Kubernetes lourde, et sans stack de monitoring consommatrice (ni Prometheus, ni Grafana).
    \item \textbf{Piliers fonctionnels à implémenter} :
    \begin{enumerate}
        \item \textbf{Canary Deployment} avec répartition pondérée du trafic ($90\%$ Stable / $10\%$ Canary).
        \item \textbf{Feature Flags (Unleash)} pour l'activation dynamique de fonctionnalités et l'injection contrôlée de pannes.
        \item \textbf{Observabilité & Health Checks} via Spring Boot Actuator.
        \item \textbf{Tests de charge automatisés} via l'outil k6.
        \item \textbf{Contrôleur de Rollback automatique} capable de basculer à chaud vers $100\%$ Stable en cas d'anomalie.
    \end{enumerate}
\end{itemize}

% ─────────────────────────────────────────────────────────────
% 2. ARCHITECTURE DE LA SOLUTION
% ─────────────────────────────────────────────────────────────
\section{Architecture Globale de la Solution}

L'architecture repose sur la séparation claire des responsabilités entre le routage réseau, les instances applicatives versionnées et le système de configuration dynamique des feature flags.

\begin{figure}[H]
\centering
\begin{tcolorbox}[colback=white, colframe=primaryblue!80, width=\textwidth, title=\textbf{Schéma Architectural du Déploiement Canary}]
\centering
\vspace{0.2cm}
\begin{verbatim}
                             Clients / Internet
                                     │
                                     ▼
                          ┌─────────────────────┐
                          │  student-nginx:80   │  (Reverse Proxy / Routeur Canary)
                          └──────┬──────┬───────┘
                                 │      │
                  90% Trafic     │      │     10% Trafic
             ┌───────────────────┘      └──────────────────┐
             ▼                                             ▼
  ┌──────────────────────┐                     ┌──────────────────────┐
  │  student-backend-v1  │                     │  student-backend-v2  │
  │   (Version Stable)   │                     │   (Version Canary)   │
  └──────────┬───────────┘                     └──────────┬───────────┘
             │                                            │
             └───────────────────┬────────────────────────┘
                                 │ Évaluation des toggles
                                 ▼
                      ┌──────────────────────┐
                      │   student-unleash    │  ◄── Dashboard Web (:4242)
                      │   (Feature Flags)    │
                      └──────────┬───────────┘
                                 │ Stockage persistant
                                 ▼
                      ┌──────────────────────┐
                      │  student-unleash-db  │  (PostgreSQL 15)
                      └──────────────────────┘
\end{verbatim}
\vspace{0.2cm}
\end{tcolorbox}
\caption{Topologie multi-conteneurs pour le Canary Testing}
\label{fig:architecture}
\end{figure}

\subsection{Description des Rôles}
\begin{enumerate}
    \item \textbf{Nginx Reverse Proxy (\texttt{student-nginx})} : Point d'entrée unique sur le port 80. Il assure la distribution pondérée des requêtes API ($90\%$ vers V1, $10\%$ vers V2), achemine les requêtes d'actifs statiques vers le frontend Angular, et expose des routes directes sans canary pour les sondes Actuator (\texttt{/v1/actuator/} et \texttt{/v2/actuator/}).
    \item \textbf{Backend V1 (\texttt{student-backend-v1})} : Version de référence en production, garantie stable. Elle dessert $90\%$ des utilisateurs.
    \item \textbf{Backend V2 (\texttt{student-backend-v2})} : Version canary candidate recevant $10\%$ du trafic. Elle embarque le filtre d'injection de pannes conditionné par Unleash.
    \item \textbf{Serveur Unleash (\texttt{student-unleash})} : Serveur de gestion de drapeaux de fonctionnalités, pilotable via API REST et interface graphique web sur le port 4242.
    \item \textbf{Base de données PostgreSQL (\texttt{student-unleash-db})} : Stockage persistant pour la configuration des feature flags.
\end{enumerate}

\screenshotplaceholder[6cm]{État d'exécution des conteneurs avec docker-compose ps}

% ─────────────────────────────────────────────────────────────
% 3. IMPLÉMENTATION DES MODIFICATIONS
% ─────────────────────────────────────────────────────────────
\section{Implémentation Détaillée des Composants}

\subsection{Évolution du Backend Spring Boot}

\subsubsection{Dépendances Maven (\texttt{pom.xml})}
Pour doter l'application de capacités d'observabilité et d'intégration avec Unleash, les modules suivants ont été intégrés :
\begin{lstlisting}[language=XML]
<!-- Spring Boot Actuator pour les sondes de sante -->
<dependency>
    <groupId>org.springframework.boot</groupId>
    <artifactId>spring-boot-starter-actuator</artifactId>
</dependency>

<!-- Client officiel Unleash pour Java -->
<dependency>
    <groupId>io.getunleash</groupId>
    <artifactId>unleash-client-java</artifactId>
    <version>9.2.6</version>
</dependency>
\end{lstlisting}

\subsubsection{Traçabilité du Trafic : \texttt{VersionHeaderFilter}}
Afin d'auditer et de mesurer précisément le routage canary depuis l'extérieur (tests k6, navigateur, scripts), un filtre HTTP de priorité maximale ajoute l'en-tête \texttt{X-Backend-Version} à chaque réponse :
\begin{lstlisting}[language=Java]
@Component
@Order(Ordered.HIGHEST_PRECEDENCE)
public class VersionHeaderFilter extends OncePerRequestFilter {
    @Value("${app.version}")
    private String appVersion;

    @Override
    protected void doFilterInternal(HttpServletRequest request,
                                    HttpServletResponse response,
                                    FilterChain filterChain)
            throws ServletException, IOException {
        response.setHeader("X-Backend-Version", appVersion);
        filterChain.doFilter(request, response);
    }
}
\end{lstlisting}

\subsubsection{Injection de Pannes et Latence : \texttt{FailureInjectionFilter}}
Ce filtre est le composant central des tests de résilience. Il applique le principe de moindre surprise :
\begin{itemize}
    \item Si l'instance est \texttt{v1}, aucune panne n'est jamais injectée.
    \item Si l'instance est \texttt{v2} et que l'URI commence par \texttt{/api/}, il interroge Unleash :
    \begin{itemize}
        \item Si le flag \texttt{simulate-failure} est activé, une réponse HTTP 500 au format JSON est immédiatement retournée.
        \item Si le flag \texttt{simulate-latency} est activé, un délai artificiel de 3000 ms est imposé avant de poursuivre la chaîne de filtres.
    \end{itemize}
\end{itemize}

\begin{lstlisting}[language=Java]
if ("v2".equals(appVersion) && request.getRequestURI().startsWith("/api/")) {
    if (unleash.isEnabled("simulate-failure")) {
        response.setStatus(500);
        response.setContentType("application/json");
        response.getWriter().write("{\"error\":\"Simulated failure (canary test)\",\"version\":\"v2\"}");
        return;
    }
    if (unleash.isEnabled("simulate-latency")) {
        Thread.sleep(3000);
    }
}
\end{lstlisting}

\subsubsection{Ouverture des Sondes Actuator (\texttt{SecurityConfig})}
Dans \texttt{SecurityConfig.java}, la chaîne de filtres de sécurité a été modifiée pour autoriser les requêtes vers Actuator sans authentification JWT :
\begin{lstlisting}[language=Java]
.requestMatchers("/api/auth/**").permitAll()
.requestMatchers("/actuator/**").permitAll() // Autorisation des sondes de sante
.requestMatchers("/h2-console/**").permitAll()
\end{lstlisting}

\screenshotplaceholder[6cm]{Reponse du endpoint Actuator Health /v1/actuator/health et /v2/actuator/health}

% ─────────────────────────────────────────────────────────────
% 4. GESTION DES FEATURE FLAGS (UNLEASH)
% ─────────────────────────────────────────────────────────────
\section{Gestion des Feature Flags avec Unleash}

\subsection{Configuration des Drapeaux}
Dans le cadre de l'expérimentation, deux drapeaux de fonctionnalités de type \textit{Release} ont été créés dans le projet \texttt{default} d'Unleash :
\begin{enumerate}
    \item \texttt{simulate-failure} : Déclenche l'erreur HTTP 500 sur le backend canary.
    \item \texttt{simulate-latency} : Déclenche une latence de 3 secondes sur le backend canary.
\end{enumerate}

\screenshotplaceholder[6.5cm]{Interface web Unleash (Port 4242) avec les drapeaux simulate-failure et simulate-latency}

\subsection{Sécurité et Jetons d'API}
Le déploiement utilise des tokens initialisés dès le démarrage du conteneur via les variables d'environnement de Docker Compose :
\begin{itemize}
    \item \textbf{Token Admin} (\texttt{*:*.unleash-insecure-admin-api-token}) : Permet le pilotage automatisé et les appels curl.
    \item \textbf{Token Client} (\texttt{default:development.unleash-insecure-api-token}) : Utilisé par le SDK Java des backends.
    \item \textbf{Identifiants Web} : Nom d'utilisateur \texttt{admin}, mot de passe par défaut \texttt{unleash4all}.
\end{itemize}

\screenshotplaceholder[6cm]{Activation du feature flag simulate-failure dans l'environnement development}

% ─────────────────────────────────────────────────────────────
% 5. TESTS DE CHARGE AVEC K6
% ─────────────────────────────────────────────────────────────
\section{Tests de Charge et Validation du Canary avec k6}

\subsection{Objectifs du Test}
Le test de charge réalisé avec \textbf{k6} vise à :
\begin{enumerate}
    \item Valider empiriquement le ratio de routage de Nginx ($90\%$ vers V1 et $10\%$ vers V2).
    \item Mesurer le temps de réponse sous charge concurrente.
    \item Confirmer le respect des seuils de qualité (SLA) : $p(95) < 2000\,\text{ms}$ et taux d'erreurs $< 10\%$.
\end{enumerate}

\subsection{Scénario de Test (\texttt{scripts/k6-load-test.js})}
Le script simule 5 utilisateurs virtuels (VUs) envoyant des requêtes d'authentification en continu pendant 30 secondes :
\begin{lstlisting}[language=JavaScript]
export const options = {
    scenarios: {
        normal_traffic: {
            executor: 'constant-vus',
            vus: 5,
            duration: '30s',
            exec: 'normalTraffic',
        },
    },
    thresholds: {
        http_req_duration: ['p(95)<2000'],
        errors: ['rate<0.1'],
    },
};
\end{lstlisting}

\subsection{Résultats Obtenus}
L'exécution sous conteneur k6 a produit les résultats chiffrés suivants :
\begin{table}[H]
\centering
\caption{Bilan de l'exécution du test de charge k6}
\begin{tabularx}{0.85\textwidth}{l X r}
\toprule
\textbf{Métrique} & \textbf{Description} & \textbf{Valeur Observée} \\
\midrule
Total Requêtes & Volume global traité en 30s & \textbf{150 requêtes} \\
Distribution V1 & Requêtes traitées par le backend stable & \textbf{90.00\%} (135 req) \\
Distribution V2 & Requêtes traitées par le backend canary & \textbf{10.00\%} (15 req) \\
Taux d'erreur & Requêtes en statut HTTP $\ge 500$ & \textbf{0.00\%} \\
Latence Moyenne & Temps de réponse moyen & \textbf{11.26 ms} \\
Latence $p(95)$ & 95e percentile de la durée & \textbf{13.96 ms} (Seuil $< 2000$ ms) \\
\bottomrule
\end{tabularx}
\end{table}

\screenshotplaceholder[7cm]{Rapport d'execution du test de charge k6 affichant la repartition 90/10 et les percentiles}

% ─────────────────────────────────────────────────────────────
% 6. ROLLBACK AUTOMATIQUE
% ─────────────────────────────────────────────────────────────
\section{Contrôleur de Rollback Automatique}

\subsection{Fonctionnement de l'Agent Watchdog}
Le script \texttt{scripts/rollback-controller.sh} agit comme un chien de garde (watchdog). Il sonde périodiquement le backend canary V2 :
\begin{enumerate}
    \item Il vérifie la sonde Actuator (\texttt{/actuator/health}).
    \item Il évalue le taux d'erreur sur un lot de requêtes tests envoyées directement à V2.
    \item Si le taux d'erreur dépasse le seuil configuré ($5\%$) pendant $N$ vérifications consécutives ($N=3$), le script déclenche le rollback.
\end{enumerate}

\subsection{Mécanisme de Bascule à Chaud}
L'opération de rollback ne nécessite aucun arrêt de conteneur, garantissant l'absence de coupure de service (\textit{zero downtime}) :
\begin{lstlisting}[language=bash]
do_rollback() {
    log "ROLLBACK TRIGGERED: Removing V2 from Nginx routing"
    # Remplacement a chaud de la configuration Nginx
    docker cp nginx/nginx-rollback.conf student-nginx:/etc/nginx/conf.d/default.conf
    # Rechargement gracieux des processus workers
    docker exec student-nginx nginx -s reload
    log "ROLLBACK COMPLETE: All traffic now routed to V1"
}
\end{lstlisting}

\screenshotplaceholder[6.5cm]{Execution du rollback-controller detectant l anomalie et appliquant le rollback}

\subsection{Validation du Rétablissement Post-Rollback}
Après déclenchement du rollback alors même que la panne simulée était maintenue active sur V2 :
\begin{itemize}
    \item 30 requêtes consécutives ont été envoyées sur le point d'entrée public de l'API.
    \item \textbf{Résultat} : $100\%$ des requêtes ont retourné \texttt{HTTP 200 OK} avec l'en-tête \texttt{X-Backend-Version: v1}.
    \item Les utilisateurs finaux ont été complètement isolés de la régression.
\end{itemize}

\screenshotplaceholder[6cm]{Verification du trafic post-rollback (100 pourcent V1 et 0 erreur utilisateur)}

% ─────────────────────────────────────────────────────────────
% 7. SCÉNARIO DE DÉMONSTRATION DE BOUT EN BOUT
% ─────────────────────────────────────────────────────────────
\section{Scénario de Démonstration de Bout en Bout}

Le script interactif \texttt{scripts/demo-shift-right.sh} regroupe l'ensemble du cycle de vie du Shift-Right Testing :
\begin{enumerate}
    \item \textbf{Étape 1} : Audit de santé initial de V1 et V2.
    \item \textbf{Étape 2} : Test de conformité du routage canary ($90/10$).
    \item \textbf{Étape 3} : Activation de l'injection de panne dans Unleash.
    \item \textbf{Étape 4} : Observation de l'apparition des erreurs HTTP 500 sur V2.
    \item \textbf{Étape 5} : Déclenchement automatique du contrôleur de rollback et rétablissement du service.
\end{enumerate}

\screenshotplaceholder[7cm]{Execution complete du script de demonstration demo-shift-right.sh}

% ─────────────────────────────────────────────────────────────
% 8. DÉPLOIEMENT SUR AWS EC2
% ─────────────────────────────────────────────────────────────
\section{Guide de Déploiement Cloud (AWS EC2)}

\subsection{Configuration de l'Instance}
Le script \texttt{deploy-ec2.sh} assure l'automatisation complète de la configuration d'une instance Ubuntu neuve :
\begin{enumerate}
    \item \textbf{Allocation d'un Swap de 2 Go} : Indispensable pour éviter les erreurs \textit{Out Of Memory (OOM)} lors des compilations Maven sur machine \texttt{t2.micro}.
    \item \textbf{Installation de Docker et Docker Compose}.
    \item \textbf{Configuration des permissions} sur les scripts.
    \item \textbf{Construction et démarrage} des 6 conteneurs.
\end{enumerate}

\subsection{Règles du Groupe de Sécurité AWS}
\begin{table}[H]
\centering
\caption{Configuration des flux réseau entrants (Security Group)}
\begin{tabular}{llll}
\toprule
\textbf{Port} & \textbf{Protocole} & \textbf{Source} & \textbf{Description} \\
\midrule
22 & TCP & Votre IP /32 & Administration sécurisée via SSH \\
80 & TCP & 0.0.0.0/0 & Accès public Application \& API Canary \\
4242 & TCP & Votre IP /32 & Tableau de bord Unleash Feature Flags \\
\bottomrule
\end{tabular}
\end{table}

\screenshotplaceholder[6cm]{Console AWS EC2 montrant l instance en cours d execution et les regles du groupe de securite}

% ─────────────────────────────────────────────────────────────
% 9. CONCLUSION & BILAN
% ─────────────────────────────────────────────────────────────
\section{Conclusion et Perspectives}

Ce projet démontre avec succès la mise en œuvre pratique du **Shift-Right Testing** dans un cadre contraint :
\begin{itemize}
    \item \textbf{Robustesse} : Le déploiement canary permet d'exposer de nouvelles versions à un sous-ensemble restreint d'utilisateurs sans impacter l'ensemble du système.
    \item \textbf{Contrôle dynamique} : L'utilisation des feature flags via Unleash permet d'activer, désactiver ou tester la résilience d'un composant en temps réel sans nécessiter de recompilation ni de redéploiement.
    \item \textbf{Automatisation de la résilience} : Le contrôleur de rollback démontre qu'un mécanisme simple et autonome de supervision permet d'annuler une mise en production défaillante en quelques secondes sans interruption de service.
    \item \textbf{Sobriété technique} : L'ensemble de la solution fonctionne sans la lourdeur d'un cluster Kubernetes ni d'une stack Prometheus/Grafana, tout en respectant scrupuleusement les contraintes d'une instance EC2 \texttt{t2.micro}.
\end{itemize}

\end{document}
